% 01_introduction.tex
% Chapter 1: Introduction

\section{Introduction}

For decades, humanity has gazed at the stars, but in the mid-20th century, we began looking back down at Earth to connect our global civilization. The concept of placing a wireless repeater in a geostationary equatorial orbit—famously envisioned by science-fiction writer and engineer Arthur C. Clarke in 1945—has transformed from a bold dream into the invisible backbone of modern telecommunications. Today, satellite links bridge continents, broadcast live events to millions, and provide crucial communications in the aftermath of natural disasters. This term project presents an in-depth, hands-on study of these remarkable space-based networks.

\subsection{Fundamental Principles of Satellite Communications}
At the heart of geostationary satellite communications lies a beautiful dance of celestial mechanics and physics. By placing a satellite in a circular orbit directly above the Earth's equator at an altitude of approximately 35,786~km (a semi-major axis of 42,164~km), the satellite's orbital velocity matches the Earth's rotation exactly. 

To an observer on the ground, the satellite appears completely frozen at a fixed longitude in the sky. This simple orbital trick has massive engineering benefits: ground stations do not need complex, continuous mechanical tracking systems. Instead, a simple dish antenna can be pointed at a single point in the sky and remain locked onto the satellite indefinitely. However, achieving this "fixed" link requires careful RF link budget calculations to overcome the immense losses accumulated over a 37,000~km journey through the vacuum of space and the Earth's atmospheric layers.

\subsection{Bent-Pipe Transponder Architecture}
The satellite payload modeled in this project is a transparent, classic \textit{bent-pipe} transponder. Much like its name suggests, a bent-pipe satellite acts as a simple "mirror in the sky." The satellite does not demodulate, decode, or process the data it receives. 

Instead, the payload is transparent: it captures the weak uplink signal arriving from Earth, filters out noise using a low-noise amplifier (LNA), translates the frequency to the downlink band using a local oscillator and mixer, amplifies the signal to high power using a Traveling-Wave Tube Amplifier (TWTA) or a Solid-State Power Amplifier (SSPA), and retransmits it back to the ground. 

While this simplicity makes the satellite incredibly flexible—independent of modulation, coding, or data formats—it introduces a major engineering challenge: uplink noise and interference are amplified and retransmitted right along with the desired signal. As a result, the uplink and downlink link budgets are physically coupled, requiring careful engineering to ensure the combined carrier-to-noise ratio ($C/N_0$) is strong enough for reliable demodulation.

\subsection{Frequency Allocation and Polarization}
This project explores a high-capacity link operating in the Ku-band, which is the commercial workhorse for Direct-to-Home (DTH) television broadcasting and VSAT business networks. Our uplink operates at a nominal frequency of 14.0~GHz, while the downlink is frequency-translated to 12.0~GHz. We use circular polarization (with a polarization tilt angle $\tau = 45^\circ$), which eliminates the need for precise mechanical polarization alignment at the earth terminals and provides excellent immunity to Faraday rotation in the ionosphere.

The Ku-band frequency band represents a sweet spot in space communications. It provides narrow, highly directional beam patterns and high antenna gains for compact dish sizes, preventing interference with adjacent satellites. However, this high frequency has a major vulnerability: water molecules. Rain, clouds, and atmospheric gases absorb and scatter Ku-band signals, presenting a dynamic propagation challenge that must be carefully calculated and mitigated to prevent complete system outages.

\subsection{Scenario Definition}
To ground our study in a realistic commercial context, we model a high-capacity, one-way transparent link between two historic capitals—Rome and Ankara—connecting Europe and Anatolia via the commercial \textbf{Eutelsat Hotbird 13G} satellite:
\begin{equation*}
    \text{GS1 (Rome)} \xrightarrow{\text{Uplink: } 14.0\text{ GHz}} \text{Eutelsat Hotbird 13G (13E)} \xrightarrow{\text{Downlink: } 12.0\text{ GHz}} \text{GS2 (Ankara)}
\end{equation*}
*   \textbf{Ground Station 1 (GS1 Rome):} Located in Rome, Italy ($41.9028^\circ$ N, $12.4964^\circ$ E, altitude 0.0~km). It acts as the broadcasting terminal (transmitting earth station), utilizing a 2.4-meter parabolic dish antenna with a 20.0~Watt high-power amplifier (HPA).
*   \textbf{GEO Satellite (Hotbird 13G):} Positioned at $13.0^\circ$ E longitude on the geostationary equatorial arc. It is modeled with a constant downlink EIRP footprint value of 50.0~dBW over the coverage zone and a receive G/T value of 3.0~dB/K at the uplink beam.
*   \textbf{Ground Station 2 (GS2 Ankara):} Located in Ankara, Turkey ($39.9334^\circ$ N, $32.8597^\circ$ E, altitude 0.0~km). It acts as the receive-only terminal, utilizing a 0.9-meter parabolic dish receiver with a receiver system noise temperature ($T_{sys}$) of 150.0~Kelvin.

\subsection{Structure of the Report}
This report provides a structured, comprehensive analysis of this geostationary link. Chapter 2 details the classical link budget equations based on Pratt's textbook to establish our open-sky baseline. Chapter 3 presents advanced optional modeling layers (such as adjacent satellite interference, intermodulation distortion, digital BER/FEC, dynamic system noise temperature, and apparent GEO orbital drift). Chapter 4 models the standards-based ITU-R P.618 atmospheric suite alongside a DVB-S2 Adaptive Coding and Modulation (ACM) selector. Chapter 5 discusses the numerical results and provides detailed physical interpretations of our generated 2D contour maps, followed by a concluding summary in Chapter 6.
